\begin{workedexample}[Worked Example: Flip angle from B1]
The flip angle produced by an RF pulse is $\alpha = \gamma\int B_1(t)\,dt$
--- proportional to the transmit field. Because $B_1$ is non-uniform (especially at
high field), the actual flip angle varies across the image, altering contrast.
$B_1$ mapping measures this: the double-angle method acquires images at nominal
$\alpha$ and $2\alpha$ and uses the signal ratio $S_{2\alpha}/S_\alpha = 2\cos\alpha$
to solve for the true local flip angle.
\end{workedexample}

\begin{keytakeaways}
\begin{itemize}[leftmargin=1.4em,itemsep=2pt]
\item Flip angle $\alpha = \gamma\int B_1\,dt$; a non-uniform $B_1$ makes $\alpha$ vary spatially.
\item $B_1$ inhomogeneity worsens at high field (shorter RF wavelength in tissue).
\item Double-angle method: $S_{2\alpha}/S_\alpha = 2\cos\alpha$ recovers the true flip angle.
\end{itemize}
\end{keytakeaways}

\begin{exercises}
\item Flip angle is proportional to which field?
\item What does $S_{2\alpha}/S_\alpha = 2\cos\alpha$ let you solve for?
\item Does $B_1$ non-uniformity get worse at higher field?
\end{exercises}

\furtherreading{Haacke~\cite{haacke}; Nishimura~\cite{nishimura}.}
